\section{Numerical experiments}\label{sec:simulations}
All experiments use $n_1=n_2=10$ and $\boldsymbol\Sigma_1=\boldsymbol\Sigma_2=\boldsymbol I_d$, so $\tau_1=\tau_2=d/10$. The first $d/2$ entries of $\boldsymbol\mu$ are equal and the remaining entries are zero, with $\lVert\boldsymbol\mu\rVert^2=\Delta$. Each outlier has covariance $\sigma_0^2\boldsymbol I_d$ and an offset supported on a separate block of the last $d/2$ coordinates, so (A2) holds. The offset norm is chosen to attain the target effective distance.

The computation solves the hard-margin dual directly. Numerically, it optimizes $\boldsymbol\gamma=d\boldsymbol\alpha$ against the normalized signed Gram matrix; this is equivalent to using a soft-margin upper bound so large that it is never active. The complete implementation, fixed seed, raw CSV files, software versions, and figure-generation code are included in \texttt{experiments/run\_experiments.py} and \texttt{results/}.

\paragraph{Why the replication counts are adequate here.}
The experiments are validation checks of deterministic asymptotic formulas, not Monte Carlo estimates of an unknown scientific effect. Each experiment compares an entire parameter grid against a prespecified curve or phase boundary, and the predicted changes are large (including exactly zero limiting inlier mass and near-unit changes in error probability). Ten replications are used for Experiments 1 and 4, twenty for the convergence-rate study, eight for Experiments 5 and 6, and six for Experiments 3 and 7. These counts are intentionally modest because one replication contains up to forty thousand coordinates and the aim is to check structural agreement rather than to report high-precision simulation estimates. Reproducibility is provided by a fixed seed, and Experiment 2 additionally reports Monte Carlo standard errors in its CSV output. The conclusions are based on agreement across grid points and across seven distinct checks, not on significance testing at a single point.

\subsection{Experiment 1: dual value and weight allocation}
Set $d=20{,}000$, $\Delta=10{,}000$, and hence $\Delta_*=14{,}000$. Add $m\in\{1,3\}$ outliers with $\sigma_0=0.2$ and choose their common distance so that $\kappa=r\Delta_*$ for $r\in\{0.1,0.3,1,3,10,30\}$. Breakdown occurs at $\kappa/\Delta_*\approx0.179$. Table~\ref{tab:exp1} and Figure~\ref{fig:exp1} compare ten-replication means with Theorem~\ref{thm:one-sided}. The two values of $m$ produce essentially the same results, as predicted by the reduction to $\kappa$.
\begin{table}[htbp]
\centering\small
\caption{Experiment 1: simulated and theoretical dual quantities, multiplied by $d$.}
\label{tab:exp1}
\resizebox{\textwidth}{!}{\begin{tabular}{rrrrrrrr}
\toprule
$m$ & $\kappa/\Delta_\ast$ & $dL$ (sim.) & $dF$ (th.) & $db$ (sim.) & $db^\ast$ (th.) & $ds$ (sim.) & $ds^\ast$ (th.) \\
\midrule
1 & 0.100 & 6.689 & 6.780 & 0.000 & 0.000 & 6.689 & 6.780 \\
1 & 0.300 & 4.786 & 4.780 & 0.922 & 0.934 & 3.864 & 3.846 \\
1 & 1.000 & 3.529 & 3.524 & 2.194 & 2.190 & 1.335 & 1.333 \\
1 & 3.000 & 3.035 & 3.090 & 2.580 & 2.625 & 0.456 & 0.465 \\
1 & 10.000 & 2.959 & 2.928 & 2.815 & 2.786 & 0.145 & 0.142 \\
1 & 30.000 & 2.964 & 2.881 & 2.912 & 2.833 & 0.051 & 0.048 \\
3 & 0.100 & 6.910 & 6.780 & 0.000 & 0.000 & 6.910 & 6.780 \\
3 & 0.300 & 4.762 & 4.780 & 0.871 & 0.934 & 3.891 & 3.846 \\
3 & 1.000 & 3.490 & 3.524 & 2.159 & 2.190 & 1.331 & 1.333 \\
3 & 3.000 & 3.101 & 3.090 & 2.636 & 2.625 & 0.465 & 0.465 \\
3 & 10.000 & 2.944 & 2.928 & 2.803 & 2.786 & 0.141 & 0.142 \\
3 & 30.000 & 2.898 & 2.881 & 2.852 & 2.833 & 0.046 & 0.048 \\
\bottomrule
\end{tabular}
}
\end{table}
\begin{figure}[htbp]
\centering
\includegraphics[width=.72\textwidth]{experiment1.pdf}
\caption{Experiment 1: dual optimum against the effective-distance ratio. The vertical dotted line is the breakdown threshold.}
\label{fig:exp1}
\end{figure}

\subsection{Experiment 2: convergence in dimension}
Vary $d$ from $1{,}000$ to $40{,}000$, set $\Delta=d/2$, $m=1$, and $\kappa=\Delta_*$. Over twenty replications, record the mean absolute relative errors of the dual value and class-2 inlier mass. Alongside Gaussian noise, use a standardized $t_5$ distribution. Figure~\ref{fig:exp2} shows decay under both distributions at the order suggested by the Chebyshev bounds in Lemma~\ref{lem:geometry}.
\begin{figure}[htbp]
\centering
\includegraphics[width=.72\textwidth]{experiment2.pdf}
\caption{Experiment 2: mean absolute relative errors against $d$ on log--log axes.}
\label{fig:exp2}
\end{figure}

\subsection{Experiment 3: critical number of outliers}
Set $d=20{,}000$, $\Delta=10{,}000$, and give every outlier common effective distance $D=2\Delta$. Then $\kappa=D/m$, and Corollary~\ref{cor:allocation} predicts collapse for $m>8$. Table~\ref{tab:exp3}, based on six replications, shows the class-2 inlier mass approaching zero at the predicted boundary.
\begin{table}[htbp]
\centering\small
\caption{Experiment 3: dependence on the number of outliers at fixed individual effective distance.}
\label{tab:exp3}
\resizebox{\textwidth}{!}{\begin{tabular}{rrrrrrrr}
\toprule
$m$ & $\kappa$ & $dL$ (sim.) & $dF$ (th.) & $db$ (sim.) & $db^\ast$ (th.) & $ds$ (sim.) & $ds^\ast$ (th.) \\
\midrule
1 & 4000.000 & 3.383 & 3.333 & 2.407 & 2.381 & 0.976 & 0.952 \\
2 & 2000.000 & 3.816 & 3.766 & 1.991 & 1.948 & 1.825 & 1.818 \\
4 & 1000.000 & 4.542 & 4.524 & 1.197 & 1.190 & 3.345 & 3.333 \\
6 & 666.667 & 5.176 & 5.165 & 0.472 & 0.549 & 4.704 & 4.615 \\
8 & 500.000 & 5.794 & 5.714 & 0.112 & 0.000 & 5.681 & 5.714 \\
10 & 400.000 & 6.333 & 6.154 & 0.000 & 0.000 & 6.333 & 6.154 \\
12 & 333.333 & 6.445 & 6.486 & 0.000 & 0.000 & 6.445 & 6.486 \\
16 & 250.000 & 7.008 & 6.957 & 0.000 & 0.000 & 7.008 & 6.957 \\
\bottomrule
\end{tabular}
}
\end{table}

\subsection{Experiment 4: discriminant function and misclassification}
Set $d=20{,}000$, $\Delta=3{,}000$, and $\tau=2{,}000$, so $\Delta_*=7{,}000$ and $\kappa_c/\Delta_*\approx0.19$. Add one outlier with $\sigma_0=0.05$, vary $\kappa/\Delta_*$ from $0.03$ to $5$, and evaluate 100 test scores per class in each of ten replications. Table~\ref{tab:exp4} and Figure~\ref{fig:exp4} show the predicted loss of class 1 near $\kappa_c$, while class 2 remains correctly classified.
\begin{table}[htbp]
\centering\small
\caption{Experiment 4: test-score means and misclassification rates.}
\label{tab:exp4}
\resizebox{\textwidth}{!}{\begin{tabular}{rrrrrrr}
\toprule
$\kappa/\Delta_\ast$ & $\bar y_1$ (sim.) & $y_1^\ast$ (th.) & $\bar y_2$ (sim.) & $y_2^\ast$ (th.) & $e(1)$ & $e(2)$ \\
\midrule
0.030 & 0.340 & 0.351 & 1.358 & 1.365 & 1.000 & 0.000 \\
0.060 & 0.258 & 0.262 & 1.214 & 1.208 & 1.000 & 0.000 \\
0.100 & 0.166 & 0.159 & 1.038 & 1.029 & 0.995 & 0.000 \\
0.150 & 0.039 & 0.059 & 0.934 & 0.916 & 0.680 & 0.000 \\
0.190 & 0.015 & 0.001 & 0.847 & 0.858 & 0.570 & 0.000 \\
0.250 & -0.060 & -0.065 & 0.807 & 0.792 & 0.180 & 0.000 \\
0.400 & -0.183 & -0.165 & 0.688 & 0.692 & 0.000 & 0.000 \\
0.700 & -0.252 & -0.259 & 0.593 & 0.598 & 0.000 & 0.000 \\
1.500 & -0.342 & -0.342 & 0.513 & 0.516 & 0.000 & 0.000 \\
5.000 & -0.402 & -0.401 & 0.450 & 0.456 & 0.000 & 0.000 \\
\bottomrule
\end{tabular}
}
\end{table}
\begin{figure}[htbp]
\centering
\includegraphics[width=.92\textwidth]{experiment4.pdf}
\caption{Experiment 4. Left: limiting and simulated discriminant values. Right: misclassification rates. The vertical dotted line is the theoretical transition.}
\label{fig:exp4}
\end{figure}

\subsection{Experiment 5: one-sided versus symmetric contamination}
Repeat Experiment 4 with eight replications, comparing the one-sided design with a second, mirror outlier in class 1 having the same effective distance. Under symmetric contamination, the discriminant remains at the clean limit even when the outliers carry most of the dual mass; Figure~\ref{fig:exp5} illustrates the cancellation in Corollary~\ref{cor:symmetric}.
\begin{table}[htbp]
\centering\scriptsize
\caption{Experiment 5: one-sided and symmetric two-sided contamination.}
\label{tab:exp5}
\resizebox{\textwidth}{!}{\begin{tabular}{rrrrrrrrr}
\toprule
$\kappa/\Delta_\ast$ & $\bar y_1$ one & $y_1^\ast$ one & $e_1$ one & $e_2$ one & $\bar y_1$ sym. & $y_1^\ast$ sym. & $e_1$ sym. & $e_2$ sym. \\
\midrule
0.030 & 0.352 & 0.351 & 1.000 & 0.000 & -0.483 & -0.429 & 0.000 & 0.040 \\
0.060 & 0.237 & 0.262 & 1.000 & 0.000 & -0.468 & -0.429 & 0.000 & 0.005 \\
0.100 & 0.155 & 0.159 & 0.985 & 0.000 & -0.437 & -0.429 & 0.000 & 0.000 \\
0.150 & 0.064 & 0.059 & 0.820 & 0.000 & -0.435 & -0.429 & 0.000 & 0.000 \\
0.190 & -0.002 & 0.001 & 0.525 & 0.000 & -0.433 & -0.429 & 0.000 & 0.000 \\
0.250 & -0.065 & -0.065 & 0.150 & 0.000 & -0.427 & -0.429 & 0.000 & 0.000 \\
0.400 & -0.150 & -0.165 & 0.005 & 0.000 & -0.424 & -0.429 & 0.000 & 0.000 \\
0.700 & -0.247 & -0.259 & 0.000 & 0.000 & -0.415 & -0.429 & 0.000 & 0.000 \\
1.500 & -0.350 & -0.342 & 0.000 & 0.000 & -0.422 & -0.429 & 0.000 & 0.000 \\
5.000 & -0.403 & -0.401 & 0.000 & 0.000 & -0.427 & -0.429 & 0.000 & 0.000 \\
\bottomrule
\end{tabular}
}
\end{table}
\begin{figure}[htbp]
\centering
\includegraphics[width=.92\textwidth]{experiment5.pdf}
\caption{Experiment 5. Left: discriminant values. Right: misclassification rates. Separate axes and reserved caption spacing prevent graphical elements from touching the caption.}
\label{fig:exp5}
\end{figure}

\subsection{Experiment 6: asymmetric two-sided contamination}
Fix the class-2 outlier at $\kappa_2/\Delta_*=0.06$ and vary $\kappa_1/\Delta_*$. Corollary~\ref{cor:misclass-two} predicts three regions: loss of class 2 for very small $\kappa_1$, a central region in which both classes are correct, and loss of class 1 for larger $\kappa_1$. Table~\ref{tab:exp6}, Figure~\ref{fig:exp6}, and the phase diagram in Figure~\ref{fig:phase} confirm this movement of the intercept.
\begin{table}[htbp]
\centering\scriptsize
\caption{Experiment 6: asymmetric two-sided contamination.}
\label{tab:exp6}
\resizebox{\textwidth}{!}{\begin{tabular}{rrrrrrrr}
\toprule
$\kappa_1/\Delta_\ast$ & $\bar y_1$ (sim.) & $y_1^\ast$ (th.) & $\bar y_2$ (sim.) & $y_2^\ast$ (th.) & $c^\ast$ & $e(1)$ & $e(2)$ \\
\midrule
0.010 & -1.100 & -1.135 & -0.234 & -0.220 & -0.678 & 0.000 & 0.820 \\
0.015 & -1.031 & -0.982 & -0.113 & -0.125 & -0.554 & 0.000 & 0.680 \\
0.020 & -0.898 & -0.881 & -0.008 & -0.024 & -0.452 & 0.000 & 0.525 \\
0.030 & -0.718 & -0.721 & 0.122 & 0.136 & -0.292 & 0.000 & 0.285 \\
0.060 & -0.412 & -0.429 & 0.411 & 0.429 & 0.000 & 0.015 & 0.000 \\
0.100 & -0.233 & -0.231 & 0.638 & 0.627 & 0.198 & 0.040 & 0.000 \\
0.150 & -0.103 & -0.099 & 0.757 & 0.758 & 0.330 & 0.235 & 0.000 \\
0.200 & -0.001 & -0.022 & 0.856 & 0.836 & 0.407 & 0.530 & 0.000 \\
0.250 & 0.036 & 0.029 & 0.878 & 0.886 & 0.458 & 0.625 & 0.000 \\
0.400 & 0.118 & 0.113 & 0.978 & 0.970 & 0.541 & 0.880 & 0.000 \\
1.000 & 0.214 & 0.199 & 1.123 & 1.099 & 0.649 & 0.995 & 0.000 \\
5.000 & 0.250 & 0.249 & 1.193 & 1.185 & 0.717 & 1.000 & 0.000 \\
\bottomrule
\end{tabular}
}
\end{table}
\begin{figure}[htbp]
\centering
\includegraphics[width=.92\textwidth]{experiment6.pdf}
\caption{Experiment 6. Left: discriminants and limiting intercept. Right: misclassification rates.}
\label{fig:exp6}
\end{figure}
\begin{figure}[htbp]
\centering
\includegraphics[width=.68\textwidth]{phase_diagram.pdf}
\caption{Phase diagram in the $(\kappa_1,\kappa_2)$ plane. Solid curves are misclassification boundaries, dashed curves are support-collapse boundaries, and colored points are Experiment 6 outcomes.}
\label{fig:phase}
\end{figure}

\subsection{Experiment 7: dual allocation in all three regimes}
Finally, set $d=20{,}000$, $\Delta=10{,}000$, and $\sigma_0=0.1$, and select ten pairs $(\kappa_1,\kappa_2)$ covering N, B1, and B2. Six-replication averages in Table~\ref{tab:exp7} agree with Theorem~\ref{thm:two-sided} and select the regime predicted by Corollary~\ref{cor:breakdown-two}.
\begin{table}[p]
\centering\scriptsize
\caption{Experiment 7: simulated and theoretical dual allocations, multiplied by $d$.}
\label{tab:exp7}
\resizebox{\textwidth}{!}{\begin{tabular}{rrrrrrrrrrrrl}
\toprule
$\kappa_1/\Delta_\ast$ & $\kappa_2/\Delta_\ast$ & $dL$ (sim.) & $dF$ (th.) & $da$ (sim.) & $da^\ast$ (th.) & $db$ (sim.) & $db^\ast$ (th.) & $ds_1$ (sim.) & $ds_1^\ast$ (th.) & $ds_2$ (sim.) & $ds_2^\ast$ (th.) & regime \\
\midrule
1.000 & 1.000 & 4.244 & 4.286 & 2.864 & 2.857 & 2.791 & 2.857 & 1.380 & 1.429 & 1.453 & 1.429 & N \\
0.300 & 0.300 & 7.613 & 7.619 & 2.816 & 2.857 & 2.873 & 2.857 & 4.798 & 4.762 & 4.741 & 4.762 & N \\
0.100 & 0.100 & 17.122 & 17.143 & 2.862 & 2.857 & 2.889 & 2.857 & 14.261 & 14.286 & 14.233 & 14.286 & N \\
3.000 & 0.100 & 7.457 & 7.320 & 6.746 & 6.612 & 0.000 & 0.000 & 0.711 & 0.708 & 7.457 & 7.320 & B2 \\
1.000 & 0.100 & 8.309 & 8.324 & 6.247 & 6.299 & 0.000 & 0.000 & 2.062 & 2.025 & 8.309 & 8.324 & B2 \\
0.300 & 0.100 & 11.224 & 11.220 & 5.284 & 5.296 & 0.538 & 0.418 & 5.939 & 5.923 & 10.686 & 10.801 & N \\
0.100 & 0.300 & 11.207 & 11.220 & 0.453 & 0.418 & 5.360 & 5.296 & 10.754 & 10.801 & 5.847 & 5.923 & N \\
0.100 & 1.000 & 8.272 & 8.324 & 0.000 & 0.000 & 6.239 & 6.299 & 8.272 & 8.324 & 2.034 & 2.025 & B1 \\
0.100 & 3.000 & 7.424 & 7.320 & 0.000 & 0.000 & 6.709 & 6.612 & 7.424 & 7.320 & 0.715 & 0.708 & B1 \\
0.050 & 0.050 & 31.431 & 31.429 & 2.556 & 2.857 & 3.060 & 2.857 & 28.876 & 28.571 & 28.371 & 28.571 & N \\
\bottomrule
\end{tabular}
}
\end{table}
